\section{数学归纳法}
\subsection{Peano公理**}
\begin{axiom}
    \,
    \begin{enumerate}
        \item \(0\)是一个自然数;
        \item 若\(n\)是自然数, 则\(n\)存在唯一的后继, 记作\(n++\);
        \item \(0\)不是任何自然数的后继, 即任意自然数\(n\)都有\(n++\neq 0\);
        \item 不同自然数的后继是不同的;
        \item 某个有关自然数的性质\(P\)满足: \(P(0)\)为真, 
        若\(P(k)\)为真, 则\(P(k++)\)也为真, 则对于任意的自然数\(n\)都有\(P(n)\)成立.
    \end{enumerate}
\end{axiom}




















\clearpage
\subsection{数学归纳法的应用*}
\begin{example}
    已知数列\(\{a_n\}\)满足\(a_{4n-3}=-1, a_{4n-1}=1, a_{2n}=a_n,\forall n\in \mathbb{N}^*\).
    设该数列的前\(n\)项和为\(S_n\), 则下列说法错误的是:
    \begin{enumerate}
        \setlength{\itemsep}{-\itemsep}
        \item[\rm{A.}] \(a_{31}=1\)
        \item[\rm{B.}] \(a_{2024}=-1\)
        \item[\rm{C.}] \(\exists T\neq 0,\forall n\in \mathbb{N}^*\)使得\(a_{n+T}=a_n\)
        \item[\rm{D.}] \(\forall n\in \mathbb{N}^*\)都有\(S_{2^n}=-2\)
    \end{enumerate}
\end{example}
\begin{jie}
    题目给出了数列\(\{a_n\}\)取值的一个分类, 依照角标除以\(4\)的余数来分类:
    奇数部分余\(1\)的取值为\(-1\), 余\(3\)的取值为\(-1\);
    而偶数部分根据\(a_{2n}=a_n\)划分到前两类中.对于任何一项都给出了
    一个具体的算法:先尽可能除以\(2\), 直到除以\(4\)余数为\(1\)或\(3\)为止.

    因此对于\(\mathrm{A.B.}\)选项, 直接计算得:
    \[a_{31}=a_{4\times 8-1}=1\]
    \[a_{2024}=a_{1012}=a_{506}=a_{253}=a_{64\times 4-3}=-1\]

    而对于\(\mathrm{C}.\)选项很容易认为是正确的. 但如果这真是一个周期数列,
    则\(T\)至少是\(4\)的倍数, 这是因为\(a_{4n-3}\)和\(a_{4n-1}\)的取值是\(-1\)和\(1\), 这部分的周期是\(4\).
    因此\[a_n=a_{2n}=a_{2n+T}=a_{n+\frac{T}{2}}\]
    即\(\frac{T}{2}\)也是此数列的周期, 重复这个操作, \(\frac{T}{4}\)也是此数列的周期,
    \(\forall k\in \mathbb{N}\), \(\frac{T}{2^k}\)也是此数列的周期.故此周期只可能是\(0\).
    故\(\mathrm{C}.\)选项错误.

    接下来我们使用数学归纳法来证明\(\mathrm{D}.\)选项.
    当\(n=1\)的时候, \(S_2=a_1+a_2=-1+(-1)=-2\), 成立.
    假设当\(n\leq k\)时命题都成立, 即\(S_{2^k}=-2\), 那么当\(n=k+1\)时:
    由假设可知\(\sum_{i=1}^{2^k}a_i=-2 \), 即证后面的部分\(\sum_{2^k+1}^{2^{k+1}}a_i=0\).
    这一共有\(2^k\)项, 其中的全部奇数项有\(2^{k-1}\)项, 可以分为两类:
    角标除以\(4\)余\(1\)和余\(3\)的部分, 它们数量相同正好正负抵消,  这部分的和为\(0\).
    接下来我们处理偶数部分, 也有\(2^{k-1}\)项的 , 通过\(a_{2n}=a_n\), 这\(2^{k-1}\)项的取值
    等于前\(2^k\)项中的后\(2^{k-1}\)项, 即
    \[\sum_{i=2^{k-1}+1}^{2^k}a_i=\sum_{i=2^{k-1}+1}^{2^k}a_{2i}\]
    根据假设, \(S_{2^{k-1}}=\sum_{i=1}^{2^{k-1}}=-2\), 因此\(\sum_{i=2^{k-1}+1}^{2^k}a_i=S_{2^k}-S_{2^{k-1}}=0\).
    这样我们就证明了\(S_{2^{k+1}}=S_{2^k}+0+0=-2\), 也就是当\(n=k+1\)时命题成立.
    综上所述, 我们证明了\(\forall n\in \mathbb{N}^*\)都有\(S_{2^n}=-2\).
  
    \qed
\end{jie}
\clearpage